% from ega4-13-fr.tex
\oldpage[IV-3]{197}
sionnel au point $x$ signifie, en vertu de (13.3.1, $a''$)), que $f$ est
dominant et que
$\dim_x(f^{-1}(f(x)))=\dim(f^{-1}(\eta))$ (où $\eta$ est le point générique de
$Y$)~; la définition (13.3.2) coïncide donc dans ce cas avec la définition
(13.2.2).

\begin{proposition}[13.3.3]
\label{IV.13.3.3-fr}
Soient $Y$ un préschéma localement noethérien, $f:X\to Y$ un morphisme
localement de type fini, $x$ un point de $X$, $X_j$ les composantes
irréductibles de $X$ (en nombre fini) contenant $x$. Pour que $f$ soit
équidimensionnel au point $x$, il faut et il suffit que, pour tout $j$,
$Y_j=\overline{f(X_j)}$ soit une composante irréductible de $Y$ et, en
désignant par $X_j$ et $Y_j$ les sous-préschémas fermés réduits de $X$ et $Y$
d'espaces sous-jacents $X_j$ et $Y_j$, par $f_j:X_j\to Y_j$ le morphisme
déduit de $f$ (\textbf{I}, 5.2.2), par $y_j$ le point générique de $Y_j$, que
$f_j$ soit équidimensionnel au point $x$ et que tous les nombres
$\dim(f_j^{-1}(y_j))$ soient égaux.
\end{proposition}

Cela résulte aussitôt de (13.3.1, $a'$)).

\begin{corollary}[13.3.4]
\label{IV.13.3.4-fr}
Avec les notations de (13.3.3), posons $y=f(x)$. Si $\mathcal O_{Y,y}$ est un
anneau universellement caténaire et si $f$ est équidimensionnel au point $x$,
on a
\[
\tag{13.3.4.1}
\label{IV.13.3.4.1-fr}
  \dim(\mathcal O_{X_j,x})=\dim(\mathcal O_{Y_j,y})+e
  -\deg.\operatorname{tr}_{k(y)}k(x)
\]
où $e$ est la valeur commune des nombres $\dim(f_j^{-1}(y_j))$.
\end{corollary}

Comme chacun des $\mathcal O_{Y_j,y}$, quotient de $\mathcal O_{Y,y}$, est un
anneau universellement caténaire (5.6.1), l'égalité est conséquence de (5.6.5).

\begin{corollary}[13.3.5]
\label{IV.13.3.5-fr}
Avec les notations de (13.3.4), supposons que $f$ soit équidimensionnel au
point $x$ et que $\mathcal O_{Y,y}$ soit un anneau universellement caténaire.
Si l'anneau $\mathcal O_{Y,y}$ est équidimensionnel, il en est de même de
$\mathcal O_{X,x}$~; la réciproque est vraie si l'image par $f$ de la réunion
des $X_j$ est dense dans un voisinage de $y$.
\end{corollary}

En effet, dire que $\mathcal O_{Y,y}$ est équidimensionnel signifie que les
nombres $\dim(\mathcal O_{Y'_i,y})$ sont égaux pour toutes les composantes
irréductibles $Y'_i$ de $Y$ contenant $y$, comme il résulte de (5.1.1.5)
appliqué au schéma local $\operatorname{Spec}(\mathcal O_{Y,y})$~; comme on a
la relation (13.3.4.1), il en résulte par le même raisonnement que
$\mathcal O_{X,x}$ est alors équidimensionnel. Inversement, si
$\mathcal O_{X,x}$ est équidimensionnel, tous les nombres
$\mathcal O_{Y_j,y}$ sont égaux par (13.3.4.1)~; on en déduit que
$\mathcal O_{Y,y}$ est équidimensionnel si les $Y_j$ sont \emph{toutes} les
composantes irréductibles de $Y$ contenant $y$, ce qui résulte de l'hypothèse
supplémentaire.

\begin{proposition}[13.3.6]
\label{IV.13.3.6-fr}
Soient $Y$ un préschéma localement noethérien, $f:X\to Y$ un morphisme
localement de type fini, $x$ un point de $X$, $y=f(x)$. Supposons que l'anneau
$\mathcal O_y$ soit équidimensionnel. Alors, si $\mathcal O_x$ est
équidimensionnel et si l'on a l'égalité
\[
\tag{13.3.6.1}
\label{IV.13.3.6.1-fr}
  \dim(\mathcal O_x)=\dim(\mathcal O_y)
  +\dim(\mathcal O_x\otimes_{\mathcal O_y}k(y))
\]
(cf. (13.2.7.1)) $f$ est équidimensionnel au point $x$, et la réciproque est
vraie si $\mathcal O_y$ est un anneau universellement caténaire.
\end{proposition}

Gardons les notations de (13.3.3)~; il résulte de (13.2.8) et de l'hypothèse
que $\mathcal O_x$ est équidimensionnel, que chacun des $Y_j$ est une
composante irréductible de $Y$ et que chacun des $f_j$ est équidimensionnel au
point $x$~; en outre (13.2.8), on a (compte tenu de (5.6.5))
\[
  \dim(\mathcal O_{X_j,x})=\dim(\mathcal O_{Y_j,y})
  +\dim_x(X_j\cap f^{-1}(y))-\deg.\operatorname{tr}_{k(y)}k(x).
\]

